\begin{abstract}
Reconstructing drivable digital humans from monocular video is a central
problem in multimedia content generation, with applications in telepresence,
virtual try-on, and interactive entertainment. Existing approaches face a
fundamental tension: feed-forward generators produce an animatable avatar
within seconds but suffer from limited geometric and appearance detail,
whereas per-video optimization methods attain high fidelity yet require
tens of minutes of offline training and cannot ingest new footage
incrementally. Neither paradigm supports the practically desirable workflow
of obtaining a usable avatar immediately and progressively improving it as
more video streams in. In this paper we propose \method{}, a framework that
unifies the two regimes through a \emph{continual-learning} formulation of
3D Gaussian Splatting (3DGS) avatar reconstruction. \method{} warm-starts
from a single frame using a feed-forward Gaussian generator and an
SMPL-based skeleton registration, yielding a drivable avatar in seconds.
It then refines the avatar continuously as a monocular video stream
arrives, treating each incoming clip as a continual-learning task and
mitigating catastrophic forgetting of previously observed viewpoints with
two complementary mechanisms: an Elastic Weight Consolidation (EWC) prior
that penalizes changes to Gaussian parameters of high Fisher-information
importance, and a viewpoint-stratified replay buffer that periodically
re-renders representative historical frames. A pose-conditioned selective
reactivation strategy further focuses capacity on under-observed body
regions. Experiments on monocular human video benchmarks demonstrate that
\method{} matches the quality of offline optimizers while supporting online
streaming consumption, and preserves earlier-view fidelity substantially
better than naive sequential fine-tuning. Code will be released to
facilitate reproducible research.
\end{abstract}

\begin{IEEEkeywords}
3D Gaussian Splatting, digital human, monocular video reconstruction, continual learning, drivable avatar, streaming reconstruction
\end{IEEEkeywords}
